\chapter{Configuration and Deployment Inputs}
\label{ch:configuration}

\chapteroverview
  {Identifies which tracked file owns each configurable contract, which
   configuration families are operational, and which override wins at runtime.}
  {Locate the owner and operational class before editing, apply the precedence
   rules to determine the effective value, then review artifact, installation,
   and dependency constraints before deployment.}

\section{Configuration Ownership Map}

\begin{longtable}{L{0.31\textwidth} L{0.27\textwidth} L{0.34\textwidth}}
\caption{Configuration files, primary consumers, and owned contracts}
\label{tab:configuration-ownership-map}\\
\toprule
\tableheader Source & Primary consumer & Owns \\
\midrule
\endfirsthead
\toprule
\tableheader Source & Primary consumer & Owns \\
\midrule
\endhead
\file{mfja_robot_control_config/config/robots.yaml} & \file{mfja_robot_control_config/launch/multi_robot_sim.launch.py} & Full-floor instances, enable flags, model/variant selection, and spawn pose \\ \rowseparator
\file{mfja_robot_control_config/config/robots_room_315_only.yaml} & Room-only and isolated launches & Focused-cell robot selection/defaults \\ \rowseparator
\file{mfja_robot_control_config/config/gripper_command_defaults.yaml} & Launch patching and gripper CLI & Per-instance jaw endpoints and default open/close percentages; fixed joint/controller mapping remains in \file{mfja_robot_control_config/launch/gripper_range_config.py} \\ \rowseparator
\file{mfja_robot_control_config/config/*.gui.config} & Gazebo GUI & Panel/layout profiles; runtime-safe and light are current profile defaults \\ \rowseparator
\file{mfja_robot_control_config/config/gazebo_params.yaml} & No current runtime consumer & Installed legacy publish-rate parameter file; current launches do not select it \\ \rowseparator
\file{mfja_robot_control_config/config/room_315_kinematics/rail_network_<side>.yaml} & Rail core/node & Directed nodes/segments, routing, switch guards, CSV references \\ \rowseparator
\file{mfja_robot_control_config/config/room_315_kinematics/raw_segments/*.csv} & Rail path loader & Calibrated sample order and geometry \\ \rowseparator
\file{mfja_robot_control_config/config/room_315_kinematics/rail_devices_<side>.yaml} & Rail device loader & Slots, sensors, stoppers, radii/default states \\ \rowseparator
\file{mfja_robot_control_config/config/room_315_vla/shuttle_identity.yaml} & Identity tracking, recording/evaluation, and TaskGoal fusion & Canonical eight identities and associated visuals \\ \rowseparator
\file{mfja_robot_control_config/config/room_315_vla/vla_supervisor.yaml} & Supervisor when explicitly selected & Primitive validation and recovery defaults \\ \rowseparator
\file{mfja_robot_control_config/config/room_315_vla/visual_state_runtime.yaml} & Visual runtime launch/node & Active V4 manifest/hash, mode, topics, presence and fusion policy \\ \rowseparator
\file{mfja_robot_control_config/config/room_315_vla/visual_observed_state_calibration.yaml} & Deterministic provider/smoke/training helpers & Geometry, block/slot mapping, thresholds and consistency policy; not the learned V4 node's runtime YAML \\ \rowseparator
\file{mfja_robot_control_config/config/room_315_vla/task_goal_understanding.yaml} & English parser/semantic adapter & Allowed language contract and optional local model settings \\ \rowseparator
\file{mfja_robot_control_config/config/room_315_vla/task_execution_runtime.yaml} & Task launch/gateway & Fail-closed authorization, topics, timeouts, limits, allowed V4 identity \\ \rowseparator
\file{mfja_robot_control_config/config/room_315_vla/task_execution_runtime_v4.yaml} & Frozen qualification workflow & Versioned V4 runtime snapshot; not automatically selected \\ \rowseparator
\file{mfja_robot_control_config/config/room_315_vla/pddl/domain_room315_runtime.pddl} & Task launch/PlanSys2 & Live executable planning domain \\ \rowseparator
\file{mfja_robot_control_config/config/room_315_vla/pddl/domain_room315.pddl} & Scenario/data tools & Broader expert/data-generation domain \\
\bottomrule
\end{longtable}

Exact files and their short roles are also listed in the Repository Source
Index in \cref{app:catalog}.

\section{Active, Offline, and Frozen Families}

The installed/current boundary is defined by
\file{mfja_robot_control_config/CMakeLists.txt}, the active launch files under
\file{mfja_robot_control_config/launch/}, and the versioned files under
\file{mfja_robot_control_config/config/room_315_vla/}.

The flat control package contains multiple research generations. Classify a
file before editing or invoking it:

\begin{longtable}{L{0.22\textwidth} L{0.35\textwidth} L{0.35\textwidth}}
\caption{Operational classification of active, offline, frozen, and historical files}
\label{tab:configuration-families}\\
\toprule
\tableheader Class & Examples & Rule \\
\midrule
\endfirsthead
\toprule
\tableheader Class & Examples & Rule \\
\midrule
\endhead
Base active runtime & robot registries, rail network/devices/CSV, GUI profiles, supervisor & Maintained for ordinary simulation; changes need live smoke tests \\ \rowseparator
Active learned runtime & \file{mfja_robot_control_config/config/room_315_vla/visual_state_runtime.yaml}, V4 visual contract/node, runtime PDDL, and \file{mfja_robot_control_config/config/room_315_vla/task_execution_runtime.yaml} & V4 only; requires exact external promoted/authorized artifacts \\ \rowseparator
Offline current tooling & V4 scenario, dataset, training, evaluation, candidate, acceptance, fault tools & May create new artifacts; never mutates the current promoted bundle \\ \rowseparator
Frozen evidence/config & V4 training/final-test JSON, versioned task YAML, sealed manifests/results & Immutable after decision; supersede with a new version \\ \rowseparator
Historical research & V3, V3R1, Experiment-A, old final-test/compatibility adapters & Preserve for provenance; not the supported active runtime \\
\bottomrule
\end{longtable}

File names alone are insufficient proof that an artifact is current. Follow
the manifest chain and exact hashes. In particular, \code{final}, \code{v4},
or \code{runtime} in a filename does not replace promotion and authorization.
Local build, runtime, and user-output classes are defined separately in
\cref{ch:maintenance}; they never become configuration authority.

\section{Precedence Rules}

Precedence and forwarding are implemented in
\file{mfja_3rd_floor_bringup/launch/room_315_floor_common.py},
\file{mfja_robot_control_config/launch/room_315_visual_state_runtime.launch.py},
\file{mfja_robot_control_config/launch/room_315_task_execution.launch.py}, and
the parameter declarations in their launched nodes.

Runtime values can come from several layers. For a given node, use this
precedence from strongest to weakest:

\begin{enumerate}
  \item an explicit launch argument that is forwarded as a ROS parameter;
  \item a parameter override in the including launch;
  \item the selected YAML parameter file;
  \item the node's declared code default; and
  \item a historical document/example, which has no runtime authority.
\end{enumerate}

Not every launch argument becomes a node parameter. Some control launch
conditions, select files, alter the environment, or cause a local cache action.
Use \code{ros2 launch ... --show-args} for the public surface, then inspect the
forwarding code and \code{ros2 param dump} for an effective running node.

\section{Robot Registry Contract}

Source files: \file{mfja_robot_control_config/config/robots.yaml},
\file{mfja_robot_control_config/config/robots_room_315_only.yaml}, and
\file{mfja_robot_control_config/launch/multi_robot_sim.launch.py}.

The full registry defines six instances:

\begin{longtable}{L{0.20\textwidth} L{0.20\textwidth} L{0.35\textwidth} L{0.15\textwidth}}
\caption{Full robot-registry instances and room-only enablement}
\label{tab:robot-registry-contract}\\
\toprule
\tableheader Name & Model & Spawn $(x,y,z,yaw)$ & Room-only \\
\midrule
\endfirsthead
\toprule
\tableheader Name & Model & Spawn $(x,y,z,yaw)$ & Room-only \\
\midrule
\endhead
\code{kuka1} & KUKA KR6 & $(-10.6366,-6.0,1.0,3.14)$ & enabled \\ \rowseparator
\code{staubli1} & St\"aubli TX2 & $(-15.1622,-6.0,1.0,1.57)$ & enabled \\ \rowseparator
\code{yaskawa_hc10_1} & HC10 & $(-15.1622,-3.0,0.62,1.57)$ & enabled \\ \rowseparator
\code{yaskawa_hc10dt_1} & HC10DT & $(-10.4366,-3.0,0.62,1.57)$ & enabled \\ \rowseparator
\code{tiago1} & TIAGo & $(-3.0,-3.0,0.0,1.57)$ & disabled \\ \rowseparator
\code{tiago_base1} & TIAGo base & $(-6.2,-3.0,0.0,1.57)$ & absent (commented example) \\
\bottomrule
\end{longtable}

When adding a robot, preserve uniqueness of instance name, Gazebo entity,
namespace, and TF prefix. Confirm model and URDF joint names align; add bridge
rules for only the capabilities actually implemented; and validate initial
joints against plugin/model limits. A registry enable flag is used only when
the launch selector is empty; \code{robots:=all} intentionally includes
disabled entries.

\section{Rail Network and Device Contract}

Source files: \file{mfja_robot_control_config/config/room_315_kinematics/rail_network_right.yaml},
\file{mfja_robot_control_config/config/room_315_kinematics/rail_network_left.yaml},
\file{mfja_robot_control_config/config/room_315_kinematics/rail_devices_right.yaml},
\file{mfja_robot_control_config/config/room_315_kinematics/rail_devices_left.yaml},
and \file{mfja_robot_control_config/scripts/room_315_kinematic_shuttle.py}.

Network YAML schema 2 defines 12 nodes, 14 directed segments, four public
switches, and explicit routing per side. A segment's \code{csv} path, direction,
start/end node, and guarded successor are one atomic contract. Device YAML
schema 1 defines four slots, 20 configured position sensors, and four stoppers
per side; derived stopper sensors are created at runtime.
The loader follows each explicit segment \code{csv} value; the current
top-level \code{source_csv_dir} metadata is not used to resolve files.

Use this rule for changes:

\begin{longtable}{L{0.25\textwidth} L{0.32\textwidth} L{0.35\textwidth}}
\caption{Rail configuration change types and mandatory validation}
\label{tab:rail-configuration-change-impact}\\
\toprule
\tableheader Intended change & Edit & Mandatory checks \\
\midrule
\endfirsthead
\toprule
\tableheader Intended change & Edit & Mandatory checks \\
\midrule
\endhead
Move rail visually as a whole & Calibration defaults/launch offsets & All device world poses, world clearance, both sides \\ \rowseparator
Correct sampled path & Referenced CSV & Monotonic index, endpoints/tangents, length drift, topology and devices \\ \rowseparator
Change connectivity & Network YAML routes/guards & All switch combinations, no-successor FALLING behavior, PDDL topology \\ \rowseparator
Move/add sensor or slot & Device YAML ratios/points & Device loader, marker, occupancy crossing, identity DZI semantics \\ \rowseparator
Rename public segment/switch & Network, defaults mapping, PDDL, configs, tests, datasets & Wire compatibility and historical adapter decision \\ \rowseparator
Change speed/headway/delay & Launch/config parameters & Multi-shuttle, stopper/switch timing, task timeouts, acceptance \\
\bottomrule
\end{longtable}

Never tune network connectivity to repair a visual coordinate offset. Use the
device-position and calibration tools, then update the explicit numeric authority.

\section{Visual and Task Parameter Precedence}

Source files: \file{mfja_robot_control_config/launch/room_315_visual_state_runtime.launch.py},
\file{mfja_robot_control_config/launch/room_315_task_execution.launch.py},
\file{mfja_robot_control_config/config/room_315_vla/visual_state_runtime.yaml},
and \file{mfja_robot_control_config/config/room_315_vla/task_execution_runtime.yaml}.

For visual runtime, the selected YAML loads first. Launch always overrides
\code{use_sim_time}, \code{device}, \code{runtime_generation},
\code{dry_run_state_fusion}, and \code{plansys2_update_enabled}. Nonempty
\code{runtime_mode}, promotion-manifest path, and manifest-hash arguments
override their YAML keys; empty values preserve YAML.

For task execution, YAML loads first. Launch always overrides simulation time,
\code{execution_enabled}, the installed runtime PDDL path, planner service, and
external-obstacle policy. \code{enable_plansys2} controls whether the planner
process is launched; it is not a gateway parameter. The planner is configured
after 1 s and automatically activated from the inactive lifecycle state.

\begin{safetybox}[Two independent enable conditions]
\code{execution_enabled:=true} is necessary but never sufficient. The gateway
also requires exact immutable authorization/promotion files and matching
schema/checkpoint hashes. Likewise, enabling the supervisor branch does not
enable the task gateway. Preserve both separations.
\end{safetybox}

\section{Host-Specific Artifact Paths}

The active tracked examples are
\file{mfja_robot_control_config/config/room_315_vla/visual_state_runtime.yaml},
\file{mfja_robot_control_config/config/room_315_vla/task_execution_runtime.yaml},
and \file{mfja_robot_control_config/config/room_315_vla/task_execution_runtime_v4.yaml};
frozen evidence paths are recorded for provenance rather than rewritten.

A Git-tracked search of \file{mfja_robot_control_config} returns 301 lines that
contain \file{/home/tiago}. They are concentrated in frozen experiment/dataset
evidence and active V4 qualification configs. Do not blindly replace them:
hashes and frozen evidence would change, and many paths intentionally record
provenance. Other repository areas, especially report evidence, may preserve
additional host paths for their own reproducibility records.

For a new machine:

\begin{enumerate}
  \item inventory which referenced bundles exist and verify every recorded hash;
  \item restore immutable bundles outside workspace \file{src/}, or create and
        qualify a new candidate;
  \item make a host-local runtime YAML or pass supported launch overrides;
  \item keep frozen evidence unchanged; and
  \item run copied-install, shadow, acceptance, fault, and authorization checks
        before enabling execution.
\end{enumerate}

\section{Install Coverage and Portability Gaps}

Source files: \file{mfja_robot_control_config/CMakeLists.txt},
\file{mfja_robot_control_config/scripts/room_315_task_goal_semantic.py},
\file{mfja_robot_control_config/scripts/room_315_task_goal_benchmark.py},
\file{mfja_robot_control_config/scripts/room_315_task_execution_config.py}, and
\file{mfja_robot_control_config/scripts/room_315_visual_observed_state_provider.py}.

The CMake install surface is explicit and is not identical to Git:

\begin{itemize}
  \item four tracked scripts are not listed for installation;
  \item \RepositoryExecutableSourceScriptCount{} of the
        \RepositoryScriptCount{} tracked source scripts carry an executable
        source bit, while CMake lists \RepositoryInstalledScriptCount{} files;
        both dimensions matter in a symlink install;
  \item \file{mfja_robot_control_config/config/room_315_kinematics/SEGMENT_CSV_MIGRATION.md}
        and two V4 coverage JSON files are
        tracked but omitted from config installation;
  \item CMake installs the complete PDDL directory, so an ignored local
        \file{mfja_robot_control_config/config/room_315_vla/pddl/domain_room315_guide_ar.html}
        can leak into a local install even
        though a clean clone cannot reproduce it; and
  \item multiple Python utilities derive config paths using
        \code{Path(__file__).resolve().parents[1]/config}. This succeeds in the
        current symlink install but can fail in a normal copied install, where
        config belongs under package share. Use \code{ament_index_python}.
\end{itemize}

The affected copied-install pattern appears in task-goal semantic/benchmark,
task-execution config, visual observed-state provider, V4 training, and related
utilities. Treat copied-install validation as a release gate until all paths
are centralized.

\section{Runtime and Tooling Dependencies Beyond Package Manifests}

Source files: \file{mfja_robot_control_config/package.xml},
\file{mfja_robot_control_config/CMakeLists.txt},
\file{mfja_robot_control_config/scripts/setup_room315_intent_model.py},
\file{mfja_robot_control_config/scripts/room_315_vla_to_lerobot.py}, and
\file{mfja_robot_control_config/launch/room_315_runtime_acceptance.launch.py}.

The manifests cover ROS, Gazebo, OpenCV, NumPy, Pillow, YAML, Torch, and
TorchVision. Selected optional workflows additionally require undeclared
\code{llama-cpp-python==0.3.16}, an alternative \code{transformers} backend,
\code{lerobot}, \code{rosbag2_py}, \code{rosidl_runtime_py}, and Python/pip for
the intent setup tool. External program assumptions include \code{ros2 bag}
and a storage plugin, \code{gz}, \code{git}, \code{nvidia-smi}, \code{pkill},
and the ROS daemon. These are environment-specific, not base simulation
requirements, but each installed workflow must state and validate them.

The runtime-acceptance launch also has the undeclared reverse dependency on
\file{mfja_3rd_floor_bringup} described in \cref{ch:packages}. Do not add that
dependency blindly: it would create a control--bringup package cycle. Refactor
the launch ownership or shared composition instead.

Conversely, Torch and TorchVision are unconditional manifest runtime
dependencies even though base simulation does not import them and the base
rosdep procedure skips their keys. This formal/operational mismatch should be
resolved by packaging the advanced learned stack separately or defining a
supported optional-dependency policy.

\section{Configuration Review Checklist}

The executable install/test surface used by this checklist is declared in
\file{mfja_robot_control_config/CMakeLists.txt},
\file{mfja_robot_control_config/package.xml}, and the registered tests under
\file{mfja_robot_control_config/test/}.

\begin{procedurebox}[Before merging a configuration change]
\begin{enumerate}
  \item Identify the owner and whether the file is active, offline, historical,
        frozen, or local output.
  \item Validate syntax and schema with the owning loader/test, not only a YAML
        or JSON parser.
  \item Search all consumers by filename, parameter key, topic, and semantic
        value; check left/right mapping separately.
  \item Confirm launch precedence and installed-file coverage in both symlink
        and copied installs where relevant.
  \item Run the smallest focused tests plus the affected end-to-end smoke path.
  \item Update hashes/versioned manifests only by the approved candidate flow.
  \item Update the launch reference, maintenance notes, and Repository Source
        Index.
\end{enumerate}
\end{procedurebox}
